\section{Experiment Set-Up}


\subsection{Research questions Definitions}

In our experiments, we aim to answer the following research questions:
\begin{itemize}
    \item \textbf{RQ1}: Can PAI-2 achieve superior results compared to baselines?
    \item \textbf{RQ2}: Does graph traversal algorithms improve PAI efficiency compared to PAI with naive flattened retriever?
    \item \textbf{RQ3}: How PAI-2`s efficiency is varying with respect to number of generating clue-queries per step of search plan?  
\end{itemize}

To choose LLM backbone for PAI-2 in our main experiments we perform a few-shot evaluation on HotpotQA dataset. We select several LLMs from the 7-9B tier: Qwen2.5 7B, Llama3.1 7B, Granite3.3 8B and Gemma2 9B. From Table~\ref{tab:llm_selection} it can be seen that best LLM by four metrics is Qwen2.5 7B. Also, to create vector representations of memory`s stored knowledge we employ combination of dense and sparse embeddings: \textit{intfloat/multilingual-e5-large}\footnote{https://huggingface.co/intfloat/multilingual-e5-large} and BM25.

\begin{table*}[ht!]
    \centering
	\renewcommand{\arraystretch}{1.3}
	%\resizebox{\textwidth}{!}{
    \begin{tabular}{|c|c|c|c|c||c|}
    \hhline{|-|----||-|}
    \multirow{2}{*}{\textbf{Method}} & \multicolumn{4}{c||}{\textbf{LLM}} & \multirow{2}{*}{Mean} \\ 
    \hhline{|~|----||~|}
     & \textbf{Qwen2.5 7B} & \textbf{Llama3.1 8B} & \textbf{Granite3.3 8B} & \textbf{Gemma2 9B} &  \\ 
     \hhline{=====::=}
    PAI-1 & \textbf{0.60} / 0.59 / 0.41 / 0.13 & 0.52 / 0.46 / 0.38 / 0.11 & 0.59 / \textbf{0.61} / 0.44 / \textbf{0.32} & 0.54 / 0.58 / \textbf{0.48} / 0.13 & 0.56 / 0.56 / 0.43 / 0.17 \\
    \hhline{|-|----||-|}
    PAI-2 & \textbf{0.70} / \textbf{0.82} / \textbf{0.44} / \textbf{0.52} & 0.64 / 0.70 / \textbf{0.44} / 0.42 & 0.54 / 0.72 / 0.39 / 0.51 & 0.65 / 0.66 / \textbf{0.44} / 0.45 & 0.63 / 0.73 / 0.43 / 0.48 \\ 
    \hhline{=====::=}
    Mean & 0.65 / 0.7 / 0.42 / 0.32 & 0.58 / 0.58 / 0.41 / 0.26 & 0.56 / 0.66 / 0.42 / 0.42 & 0.6 / 0.62 / 0.46 / 0.29 & 0.6 / 0.64 / 0.43 / 0.32 \\ 
    \hhline{|-|----||-|}
    \end{tabular}%}
    \renewcommand{\arraystretch}{1}
    \caption{Best performance in a few-shot ablation experiment for PAI-1 and PAI-2 on HotpotQA dataset across four LLMs. Cells contain Context Relevance, Faithfulness, LLM-as-a-Judge and Groundedness scores, respectively, to identify optimal LLM, that should be used in main experiments.}
	\label{tab:llm_selection}
\end{table*}

For knowledge graph traversal in PAI-2 we select two combinations of BeamSearch (BS), WaterCircles (WC) and NaiveRetirever (NR) algorithms, presented in PAI-1~\cite{11479299}, as they give superior and comparative performance by our previous research: "BS + WC" and "BS + NR". The values of hyperparameters for the base algorithms are fixed (see Appendix~\ref{app:retriev_hyper}). During graph traversal we did not apply constraints on vertex types, but during filtering stage episodic triples are discarded.

\subsection{Summary of evaluated configurations}

Each PAI-2 configuration was evaluated on 100 question-answer pairs from each benchmark. The same LLM was used for both: responses generation using given QA configuration and corresponding memory graph construction. Consequently, for each dataset 15 distinct QA configurations were derived. In total, 90 QA configurations were evaluated; plus 44 configurations for LLM few-shot ablation study on HotpotQA. 

\subsection{Implementation details}

Our memory graph implementation consists of two main parts: a graph part and a vector part. The graph part stores textual representations of object, thesis and episodic vertices, together with their properties and relationships (edges). The Neo4j is used for this part of the system. The vector part of memory stores vector representations (embeddings) of elements from the graph part to measure semantic similarity of texts during QA pipeline execution. The Qdrant and OpenSearch are used for this part of the system to store dense and spare embeddings respectively. PAI also implements a caching mechanism for storing intermediate results of QA pipeline steps to reduce overall time, that is required to process incoming questions. It utilizes two non-relational databases: Redis and MongoDB. During our experiments, cache was enabled. All databases were hosted and run on a single machine in separate Docker containers. For our needs medium-sized LLMs (7-14B) were hosted in local Ollama Docker container. LLM inference during memory construction and QA pipeline execution was performed on a single NVIDIA TITAN RTX 24GB GPU. 

For PAI-2 evaluation, we constructed six memory graphs based on six selected/preprocessed benchmarks/datasets. An average speed of adding documents (with 492 average length) to memory per minute is approximately 1.63. Detailed characteristics of constructed memory graphs can be found in Appendix~\ref{app:constr_graphs}. It is important to note that we disable the query preprocessing stage because it does not give sufficient boost to QA pipeline accuracy (based on our ablation; LLM prompts need to be tuned additionally) and also this functionality is outside the scope of stated research questions.